\Needspace{5\baselineskip}\textbf{Problems}

\begin{enumerate}[label=\thesection.\arabic*.,ref=\thesection.\arabic*,start=39,series=book-problems]

\item\label{ex:bounded-compact-operator-products} 设 $A\in\mathcal{L}(X)$ 为紧算子, $B\in\mathcal{L}(X)$ 为有界算子. 证明 $AB$ 和 $BA$ 为紧算子.

\item\label{ex:finite-rank-spectrum-alternative-proof} 证明 Lemma~\ref{lem:finite-rank-spectrum-alternative}. 使用线性代数中的适当结果.

\item\label{ex:compact-nonzero-eigenspace-finite-dimensional} 设 $A\in\mathcal{L}(X)$ 为紧算子. 证明, 当 $\lambda\ne0$ 时, $A\phi=\lambda\phi$ 的解空间有限维.

\item\label{ex:compact-operator-singular-system-expansion} 设 $A\in\mathcal{L}(H)$ 为紧算子. 证明存在标准正交集 $\{\psi_i\}_{i=1}^N$ 和 $\{\phi_i\}_{i=1}^N$, 以及正实数 $\{\lambda_i\}_{i=1}^N$ (这里 $N$ 可以有限, 也可以无穷), 使得
\begin{equation}\label{eq:compact-operator-singular-system-problem}
A(u)=\sum_{i=1}^N\lambda_i(\psi_i,u)\phi_i.
\end{equation}
Hint: $A^*A$ 为紧自伴随算子.

\item\label{ex:compact-operator-norm-preserving-extension} 证明, 若用 Theorem~\ref{thm:bounded-operator-norm-preserving-extension} 的证明中情形 1 和情形 2 所定义的方法, 扩张从 $X$ 到 $Y$ 的紧算子 $(\mathcal{D}(A),A)$, 则扩张后仍为紧算子.

\item\label{ex:compact-adjoint-hilbert-proof} 设 $H$ 为 Hilbert 空间. 在 $X=Y=H$ 的情形下证明 Theorem~\ref{thm:compact-operator-adjoint-compact}. Hint: 使用 Theorem~\ref{thm:compact-hilbert-finite-rank-approximation}.

\item\label{ex:relative-compactness-arbitrarily-small-relative-bound} 若 $\mathcal{D}(B)\supseteq\mathcal{D}(A)$, 且只要 $x_n\in\mathcal{D}(A)$ 及 $\|Ax_n\|_Y+\|x_n\|_X$ 有界, $Bx_n$ 就有收敛子序列, 则称 $B$ 相对于 $A$ 紧. 假设 $A$ 是闭的, $B$ 是可闭的, 且 $B$ 相对于 $A$ 紧. 证明 $B$ 相对于 $A$ 有界, 并且 Problem~\ref{ex:closed-unbounded-resolvent-analyticity} 中的常数 $a$ 可以任意小. Hint: 尝试仿照 Ehrling 引理 (Lemma~\ref{thm:ehrling-compact-intermediate-norm}) 的证明.

\item\label{ex:riesz-subspace-separation} 证明以下由 F. Riesz 得到的结果. 设 $S_1$ 和 $S_2$ 为某赋范线性空间的子空间. 假设 $S_1$ 是闭的, 且 $S_1$ 为 $S_2$ 的真子集. 那么对每个 $\theta\in(0,1)$, 都存在 $x\in S_2$, 使得 $\|x\|=1$ 且
\begin{equation}\label{eq:riesz-subspace-separation-problem}
\|x-y\|\geq\theta\quad\text{for all }y\in S_1.
\end{equation}
Hint: 取 $S_2\ni w\notin S_1$, 并令 $d=\operatorname{dist}(w,S_1)$. 证明存在 $v\in S_1$, 使得
\begin{equation}\label{eq:riesz-subspace-separation-distance-hint}
d\leq\|w-v\|\leq\dfrac{d}{\theta}.
\end{equation}


\item\label{ex:infinite-dimensional-identity-noncompact} 证明, 赋范空间 $X$ 中的单位球是紧的, 当且仅当 $X$ 有限维. Hint: 使用 Problem~\ref{ex:riesz-subspace-separation}.

\end{enumerate}
